\par\Needspace{6\baselineskip}\medskip
\SourceLeaf{195}\begin{center}\hypertarget{native-195-article144}{}\pdfbookmark[2]{144.}{article-144}144.\end{center}
Sind die Bögen $\zeta,\zeta''$ gefunden, so ist die Lage der Punkte $C,C''$ gegeben, und der Abstand $CC''=2f'$ lässt sich aus $\zeta,\zeta''$ und $\epsilon'$ bestimmen. Es seien $u,u''$ die Neigungen der Großkreise $AB,A''B''$ gegen den Großkreis $CC''$ (in Abb. 4 sind dies der Reihe nach die Winkel $C''CD'$ und $180^{0}-CC''D'$). Dann gelten die folgenden Gleichungen, die den Gleichungen 3—6 aus Artikel 137 genau entsprechen.
\[
\begin{aligned}
\sin f'\sin\tfrac12(u''+u)&=\sin\tfrac12\epsilon'\sin\tfrac12(\zeta+\zeta'')\\
\sin f'\cos\tfrac12(u''+u)&=\cos\tfrac12\epsilon'\sin\tfrac12(\zeta-\zeta'')\\
\cos f'\sin\tfrac12(u''-u)&=\sin\tfrac12\epsilon'\cos\tfrac12(\zeta+\zeta'')\\
\cos f'\cos\tfrac12(u''-u)&=\cos\tfrac12\epsilon'\cos\tfrac12(\zeta-\zeta'').
\end{aligned}
\]
Die ersten beiden liefern $\tfrac12(u''+u)$ und $\sin f'$, die letzten beiden $\tfrac12(u''-u)$ und $\cos f'$. Aus $\sin f'$ und $\cos f'$ erhält man $f'$. Die Winkel $\tfrac12(u''+u)$ und $\tfrac12(u''-u)$ werden erst bei der letzten Hypothese zur Bestimmung der Lage der Bahnebene herangezogen und können bei den ersten Hypothesen außer Betracht bleiben.

Ganz entsprechend ließe sich $f$ aus $\epsilon,C'D$ und $C''D$ ableiten sowie $f''$ aus $\epsilon'',CD'',C'D''$. Für diesen Zweck sind jedoch die folgenden Formeln wesentlich bequemer.
\[
\begin{aligned}
\sin2f&=r\sin2f'\cdot\frac{n}{n'r'}\\
\sin2f''&=r''\sin2f'\cdot\frac{n''}{n'r'}\mathpunct{\raisebox{.45ex}{,}}
\end{aligned}
\]
Dabei liegen die Logarithmen der Größen $\dfrac{n}{n'r'}\mathpunct{\raisebox{.45ex}{,}}$ $\dfrac{n''}{n'r'}$ bereits aus den vorigen Rechnungen vor. Schließlich lässt sich die gesamte Rechnung noch einmal dadurch prüfen, dass $2f+2f''=2f'$ gelten muss. Ergibt sich dennoch eine Differenz, so ist sie gewiss ohne Belang, wenn alle Rechnungen möglichst genau ausgeführt wurden. Mitunter kann sie jedoch, selbst wenn durchgehend mit sieben Dezimalstellen gerechnet wurde, einige Zehntel einer Bogensekunde erreichen. Scheint es der Mühe wert, so lässt sie sich sehr leicht so auf $2f,2f''$ verteilen, dass die Logarithmen der Sinuswerte um denselben Betrag erhöht oder vermindert werden. Damit ist die Gleichung $P=\dfrac{r\sin2f''}{r''\sin2f}=\dfrac{n''}{n}$ mit der vollen von den Tafeln erlaubten Genauigkeit erfüllt. Unterscheiden sich $f$ und $f''$ nur wenig, so genügt es, die Differenz gleichmäßig auf $2f$ und $2f''$ zu verteilen.
